\section{Frontier--Closure Interaction}
\label{supp:frontier-closure-interaction}\label{app:interaction-proof}

The main paper's compressed state has two economically different margins.  A
higher frontier improves current operation and makes further operational
improvement harder.  A richer closure expands the menu of policies that
research can generate.  This section gives the complete finite interaction
theory, its primitive sufficient conditions, and the exact mechanisms that
make the sign ambiguous outside those conditions.

The results use the main paper's finite-horizon unified action value,
including cost, positive duration, operation during research, the joint
completion law, compressed updating, and lower-horizon optimized
continuation.  They add no path--outcome independence assumption.

\subsection{Closure increments and relative saturation}

\begin{definition}[Closure increment and interaction]
\label{def:frontier-closure-interaction}
For \(F_0\le F_1\) and \(C_0\subseteq C_1\), define
\[
 \Delta_CV_h(F;C_1,C_0):=V_h(F,C_1)-V_h(F,C_0)
\]
and
\[
 J_h:=\Delta_CV_h(F_1;C_1,C_0)
      -\Delta_CV_h(F_0;C_1,C_0).
\]
The two margins are substitutes on the rectangle when \(J_h\le0\) and
complements when \(J_h\ge0\).
\end{definition}

Frontier independence holds project availability, initiation cost, duration,
operation during research, and the complete joint completion law fixed across
the two frontiers.  Closure expansion makes every closure-poor action
available under the rich closure.  These restrictions make each fixed
opportunity saturate, but they do not compare saturation across projects.
The theorem therefore needs a relative condition at the action level.

\begin{theorem}[Frontier--closure substitution under relative saturation]
\label{thm:frontier-closure-substitution}
Consider the four realizable compressed states
\[
 (F_0,C_0),\quad(F_0,C_1),\quad(F_1,C_0),\quad(F_1,C_1),
 \qquad F_0\le F_1,\quad C_0\subseteq C_1.
\]
Suppose research primitives are frontier independent and the closure-rich
menu contains the closure-poor menu.  Suppose also that, at every finite
Bellman node, every action \(a_1\) feasible with \(C_1\) and every action
\(a_0\) feasible with \(C_0\) satisfy
\begin{equation}\label{eq:t7-relative-saturation}
 Q_h(F_1,C_1,a_1)-Q_h(F_1,C_0,a_0)
 \le
 Q_h(F_0,C_1,a_1)-Q_h(F_0,C_0,a_0).
\end{equation}
Here \(Q_h\) is the unified action value, including cost, positive duration,
operation during research, the joint completion law, compressed updating,
and lower-horizon optimized continuation.  Then, for every finite horizon and
belief,
\[
 \Delta_CV_h(F_1;C_1,C_0)
 \le
 \Delta_CV_h(F_0;C_1,C_0),
 \qquad\text{equivalently }J_h\le0.
\]
\end{theorem}

At a fixed finite horizon and belief, for the four realizable corners choose
\(a_{11}\) maximizing at \((F_1,C_1)\) and \(a_{00}\) maximizing at
\((F_0,C_0)\).  Finiteness of the action menus supplies these maximizers.
Frontier independence transports them to the comparison corners; no
probabilistic independence is asserted.  Relative action saturation and
finite maximum attainment give
\[
\begin{aligned}
 V_h(F_1,C_1)-V_h(F_1,C_0)
 &\le Q_h(F_1,C_1,a_{11})-Q_h(F_1,C_0,a_{00})\\
 &\le Q_h(F_0,C_1,a_{11})-Q_h(F_0,C_0,a_{00})\\
 &\le V_h(F_0,C_1)-V_h(F_0,C_0).
\end{aligned}
\]
This is exactly \(J_h\le0\).  Neither detectability nor contraction is used.
Without the middle inequality, individual project saturation does not control
optimizer switching.

\subsection{Primitive common-gap sufficient conditions}

\begin{proposition}[Primitive common-gap frontier saturation]
\label{prop:primitive-frontier-saturation}
Maintain the realizable rectangle, frontier-independent primitives, and menu
inclusion of \cref{thm:frontier-closure-substitution}.  Suppose that at every
finite Bellman node and belief the four action-value tables admit
\begin{equation}\label{eq:primitive-common-gap}
\begin{aligned}
 Q_h(F_i,C_1,a)
   &=B_{h,i}+\eta^1_{h,a}+\lambda^1_{h,a}g_{h,i},\\
 Q_h(F_i,C_0,a)
   &=B_{h,i}+\eta^0_{h,a}+\lambda^0_{h,a}g_{h,i},
 \qquad i\in\{0,1\},
\end{aligned}
\end{equation}
where \(g_{h,1}\le g_{h,0}\), every rich-menu exposure satisfies
\(\lambda^1_{h,a}\ge0\), and every feasible poor-menu action satisfies
\(\lambda^0_{h,a}=0\).  Then relative action saturation
\eqref{eq:t7-relative-saturation} holds.  Consequently
\[
 \Delta_CV_h(F_1;C_1,C_0)
 \le
 \Delta_CV_h(F_0;C_1,C_0)
\]
for every finite horizon and belief.
\end{proposition}

The subclass is economically narrow but primitive.  A fixed descendant
profile supplies the common positive-part gap \(g\); frontier-independent
cost, duration, operation, candidate, admission, and fixed-continuation terms
enter the intercepts and nonnegative exposures.  The poor closure is
Continue-only in the leading case; more generally its feasible actions may
have fixed payoffs but no exposure to the descendant gap.  The shared
decomposition across frontiers is the required fixed-continuation
factorization.  It is an assumption defining this subclass, not a claim that
arbitrary unified continuation preserves the form.

There is a transition-level sufficient condition for the displayed gap order.
Let the low- and high-frontier gaps share the process belief kernel \(P\) and
discount \(\beta\ge0\), and let
\[
\begin{aligned}
 g_{0,i}(b)&=\bar g_i(b),\\
 g_{h+1,i}(b)
   &=s_i(b)+\beta\sum_{b'}P(b,b')g_{h,i}(b'),
 \qquad i\in\{0,1\}.
\end{aligned}
\]
If \(\bar g_1\le\bar g_0\) and \(s_1\le s_0\) pointwise, positivity of the
fixed transition operator gives \(g_{h,1}\le g_{h,0}\) by finite-horizon
induction.  Thus the common-gap proposition is recursion stable when the
descendant continuation uses this fixed kernel and remains outside the
action-dependent optimizer.

The induction makes this preservation step explicit.  Assume
\(\beta\ge0\), \(P(b,b')\ge0\), \(\bar g_1(b)\le\bar g_0(b)\), and
\(s_1(b)\le s_0(b)\) pointwise.  The terminal order is the induction base.
If \(g_{h,1}(b')\le g_{h,0}(b')\) for every \(b'\), then
\[
\begin{aligned}
 g_{h+1,1}(b)-g_{h+1,0}(b)
   ={}&s_1(b)-s_0(b)\\
     &+\beta\sum_{b'}P(b,b')
       \bigl(g_{h,1}(b')-g_{h,0}(b')\bigr)
   \le0.
\end{aligned}
\]
Hence \(g_{h,1}\le g_{h,0}\) for every finite \(h\).  Row mass one gives the
kernel its probability interpretation; positivity, rather than normalization,
is the sign-bearing property in the induction.

At each node, fixed project and continuation primitives give
\[
 Q_h(F_i,C_j,a)=B_{h,i}+\eta^j_{h,a}
                  +\lambda^j_{h,a}g_{h,i}.
\]
Take a rich action \(a_1\) and a feasible poor action \(a_0\).  The
frontier-specific base cancels within their relative payoff.  Because every
poor action has \(\lambda^0_{h,a_0}=0\), its change from the low to the high
frontier is
\[
\begin{aligned}
 &\bigl[Q_h(F_1,C_1,a_1)-Q_h(F_1,C_0,a_0)\bigr]\\
 &\quad-
 \bigl[Q_h(F_0,C_1,a_1)-Q_h(F_0,C_0,a_0)\bigr]
 =\lambda^1_{h,a_1}(g_{h,1}-g_{h,0})\le0,
\end{aligned}
\]
using \(\lambda^1_{h,a_1}\ge0\).  This is precisely relative action
saturation.  Frontier independence transports the high-rich and low-poor
maximizers to the other corners, and the finite-max comparison above then
yields \(J_h\le0\).

\subsection{Complementarity and sign-ambiguous cases}

The zero-exposure restriction cannot be replaced by merely ordering the added
project's exposure above the incumbent project's.  Because Continue is always
available and has zero exposure, its advantage over a positive-exposure poor
project rises when the common gap shrinks.  An exact 2,430-row rational search
finds 648 failures of all-pairs relative saturation under the broader exposure
order and none among the 810 zero-poor-exposure rows.  This is finite
falsification and validation, not proof; the proposition is separately kernel
checked.

The substitution theorem is sharp.  Exact rational finite examples attain
all three interaction signs:
\[
\begin{array}{lccc}
\toprule
\text{construction}
  &\Delta_CV(F_0)&\Delta_CV(F_1)&J\\
\midrule
\text{fixed opportunity with saturation}&2&1&-1\\
\text{frontier-dependent success }0\to1&0&1/2&1/2\\
\text{separable closure bonus}&3/2&3/2&0\\
\bottomrule
\end{array}
\]

\begin{proposition}[Complementarity from project switching]
\label{prop:t7-interaction-examples}
Frontier independence and individual project saturation do not imply
substitution.  Let \(\beta=1/2\), the descendant payoff be \(10\), and
\((F_0,F_1)=(0,8)\).  The old project has fixed
\((\pi,\kappa)=(1,2)\), while closure adds a project with fixed
\((\pi,\kappa)=(1/2,0)\).  Their premia fall respectively
\[
  3\to0
  \qquad\text{and}\qquad
  5/2\to1/2.
\]
Nevertheless, optimization gives
\[
  \Delta_CV(F_0)=0,\qquad
  \Delta_CV(F_1)=1/2,\qquad
  J=1/2.
\]
\end{proposition}

The added project's advantage over the old project rises from
\(-1/2\) to \(1/2\), violating relative action saturation.  This produces
strict complementarity even though both project premia decline individually.
The example remains outside
\cref{prop:primitive-frontier-saturation}: its poor-menu project has positive
common-gap exposure.  Frontier-dependent success provides a second exact
complementarity mechanism; the fixed-opportunity and separable examples show
that substitution and zero interaction are also attainable.  Hence neither
closure enrichment nor frontier improvement has an unconditional interaction
sign.

The numerical interaction response remains evidence separate from the
theory.  It preserves the registered 2,430-row frontier--project boundary grid
and five canonical four-corner fixtures: primitive strict substitution, the
saturation boundary, optimizer-switching complementarity,
frontier-dependent-success complementarity, and nonconstant separability.
Every row reports the two closure increments, their cross difference, the
four selected projects, four realizability flags, and the primitive
sufficient-condition certificate.  Only jointly realizable four-corner rows
enter any sign frequency; nonrealizable diagnostic rows remain recorded but
are not aggregated.
